\recipe{Tempest’s family pudding (Jam Roly Poly)}
\label{recipe:jam_roly_poly}
\makes{6 portions}
\worktime{\SI{15} \minute}
\totaltime{\SI{1} \hour}
\source{Tempest}


\begin{ingredients}

  \item 50g salted butter (cold and cut into chunks, plus extra for greasing)
  \item 250g self-raising flour (plus extra for rolling)
  \item 1 vanilla pod (seeds scraped out)
  \item 50g shredded suet
  \item 150ml milk
  \item 100g/4oz raspberry or plum jam, or a mixture
  \item Alpro vanilla yogurt
\end{ingredients}


\begin{method}
  	Put a deep roasting tin onto the bottom shelf of the oven, and make sure that
there’s another shelf directly above it. Pull the roasting tin out on its shelf,
fill two-thirds with boiling water from the kettle, then carefully slide it back
in. Heat oven to 180C/160C fan/gas 4. Tear off a large sheet of foil and
greaseproof paper (about 30 x 40cm). Sit the greaseproof on top of the foil
and butter it.

Tip butter, flour and vanilla seeds into a food processor; pulse until the
butter has disappeared. Tip into a mixing bowl. Stir through the suet, pour in
the milk and work together with a cutlery knife until you get a sticky dough.
You may need a drop more milk, depending on your flour.

Tip the dough out onto a floured surface, quickly pat together to smooth,
then roll out to a square roughly 25 x 25cm. Spread the jam all over, leaving
a gap along one edge, then roll up from the opposite edge. Pinch the jam
free edge into the dough where it meets, and pinch the ends roughly, too.
Lift onto the greased paper, join-side down, loosely bring up the paper and
foil around it, then scrunch together along the edges and ends to seal. The
roly-poly will puff quite a bit during cooking so don’t wrap it tightly. Lift
the parcel directly onto the rack above the tin and cook for 1 hr.

Let the pudding sit for 5 mins before unwrapping, then carefully open the
foil and paper, and thickly slice to serve.

  
\end{method}


\begin{figure}[h]
  	\centering
  	
  	\begin{tikzpicture}
  		% Save the image in a box to get its dimensions
  		\sbox0{\includegraphics[width = 0.6\linewidth]{images/MissingImage.png}}%
  
  		% Draw a clipping rectangle with rounded corners
  		\path[clip, rounded corners=0.5cm] (0,0) rectangle (\wd0,\ht0);
  
  		% Place the image inside the clipped area
  		\node[inner sep=0pt, outer sep=0pt, anchor=south west] at (0,0) {\usebox0};
	\end{tikzpicture}

  \caption*{caption for image with date and information, 1st January 2026}
\end{figure}

\begin{notes}
Sit on the floor of a campsite in a
circle and eat from 1 bowl for extra
atmosphere.
\end{notes}

